% Generated draft: requires mathematical and source review.
\documentclass[11pt]{article}
\usepackage[margin=25mm]{geometry}
\usepackage{fontspec}
\setmainfont{Latin Modern Roman}
\usepackage{amsmath,amssymb,mathtools,unicode-math}
\setmathfont{Latin Modern Math}
\usepackage{enumitem,needspace,xurl,hyperref}
\setlength{\parindent}{0pt}
\setlength{\parskip}{5pt}
\setlength{\emergencystretch}{3em}
\newcommand{\N}{\mathbb{N}}
\newcommand{\Z}{\mathbb{Z}}
\newcommand{\Q}{\mathbb{Q}}
\newcommand{\R}{\mathbb{R}}
\newcommand{\C}{\mathbb{C}}
\newcommand{\F}{\mathbb{F}}
\newcommand{\A}{\mathbb{A}}
\newcommand{\PP}{\mathbb{P}}
\newcommand{\E}{\mathbb{E}}
\newcommand{\eps}{\varepsilon}
\newcommand{\dd}{\,\mathrm{d}}
\begin{document}
\section*{UnsolvedMath ID: 6700054}
\textbf{Source problem code:} catalogue entry
\subsection{Definitions and mathematical statement}
Let $X=\Sigma^{n-1}\times[-1,1]$ be the smooth product band, endowed with a Riemannian metric $g$. The product notation describes the topology, not necessarily a product metric. The stated curvature hypothesis is
\[
\operatorname{sec}_g(\Pi)=1
\]
for every tangent two-plane $\Pi$; it is not a scalar-curvature-only assumption. A literal Riemannian product metric would have mixed sectional curvatures zero and would not satisfy this condition.

The source observes that there is ``no apparent non-trivial bound'' on the width even under this constant-sectional-curvature hypothesis. It does not define width or state a numerical conjecture. If ordinary band width is intended, write
\[
\operatorname{width}_g(X)
=\inf\{d_g(p,q):p\in\Sigma\times\{-1\},\ q\in\Sigma\times\{1\}\}.
\]
A precise investigation under that convention would ask for the supremum of this quantity over admissible curvature-one metrics and whether a bound depends only on dimension or also on $\Sigma$ and boundary data. The source does not fix those quantifiers, compactness/completeness assumptions, or even whether this is the intended notion of width. Its absence of an apparent bound is not asserted here as a proved unboundedness theorem.

The supplied research summary speculates about scalar-curvature limitations, but that discussion does not replace the actual sectional-curvature-equals-one hypothesis.
\subsection{Short English statement}
Investigate width bounds for a product-topology band with sectional curvature exactly one. The extract states that no bound is apparent but does not define width or give the manifold and metric hypotheses needed for a precise extremal question.
\subsection{Sources}
\begin{itemize}
\item Misha Gromov, \emph{101 Questions, Problems and Conjectures around Scalar Curvature} (2017), printed tag [?58], occurrence 57, PDF page 59, \url{https://www.ihes.fr/~gromov/wp-content/uploads/2018/08/101-problemsOct1-2017.pdf}; accessed 2026-07-29, PDF-text-normalized extraction. The provenance describes the public author PDF as incomplete.
\item Local registry \texttt{data/unsolvedmath/problems.json}, permanent UnsolvedMath ID $6700054$, source code catalogue entry, \emph{Scalar Curvature Question [?58]}. Dataset attribution: Ulam AI and the UnsolvedMath Contributors.
\end{itemize}
\end{document}
